\section{Methodology}
\label{sec:methodology}

\subsection{Protocol}

One trial has two phases:

\begin{enumerate}
\item \textbf{Warm-up} (\ResWarmupDays{} simulated days). Nominal operations with
the engine in learning mode. Behavioural rules build their profiles and emit
nothing.
\item \textbf{Measurement} (\ResMeasureDays{} simulated days,
\ResEventsPerWindow{} events). Nominal operations continue while the five attacks
are injected at staggered start times, separated by \SI{20}{\hour} plus a
per-trial random jitter of up to \SI{8}{\hour}. The spacing makes every alert
attributable to at most one attack; the jitter prevents each trial from placing
attacks at the same point of the orbital cycle, which would report schedule
artefacts as variance.
\end{enumerate}

We run \ResTrials{} independent trials with different seeds and report means with
\SI{95}{\percent} normal-approximation confidence intervals.

\subsection{The benign workload}

A testbed with no benign anomalies reports a false-positive rate of zero and says
nothing useful. Ours generates four classes of benign noise, chosen because each
is a real analyst's daily experience: operators mistyping passwords (occasionally
enough times in a row to resemble a spray); operators connecting from mobile or
travel networks; the radio link corrupting frames at low elevation; and a
legitimate weekly reprocessing job that bulk-reads the archive in a pattern
superficially identical to exfiltration. Their \emph{rates} are chosen, not
measured from a real mission --- a limitation we return to in
\S\ref{sec:limitations}.

\subsection{Scoring definitions}

\begin{description}
\item[True positive.] An alert at least one of whose evidence events carries a
non-benign ground-truth label. Attribution therefore follows causality, not
temporal coincidence: an alert that merely fires during an attack window, caused
by unrelated benign activity, counts as a false positive.
\item[Time to detect (TTD).] The first attributed alert's actionable timestamp
minus the attack's first malicious action. It includes the pipeline latency of
the tier that produced the alert. For the replay scenario the reference point is
the first re-transmission, not the passive radio capture, because capture leaves
no observable in the ground segment.
\item[Time to contain (TTC).] The first automated containment action on the
attack's entity minus the same reference point.
\item[Impact.] Malicious actions that took effect --- commands executed, objects
read, forged frames accepted --- reported both in total and, separately, as
critical commands executed and megabytes exfiltrated, because ``one executed
thruster firing'' and ``one object read'' are not comparable harms.
\item[Impact after containment.] The subset of the above occurring after the
first containment action: what the response did \emph{not} prevent.
\end{description}

\subsection{Deduplication as part of the system under test}

Without a per-entity cooldown, a single incident produces an alert storm and the
false-positive rate becomes meaningless. The cooldown (\SI{300}{\second} by
default) is therefore an explicit, swept parameter rather than a hidden
implementation detail.

\subsection{Ablation and sensitivity}

The ablation repeats the whole campaign with one control disabled at a time
across seven preventive and responsive mechanisms, plus an all-disabled
configuration. The sensitivity analysis sweeps four tuning parameters --- the
frame-integrity threshold, the batch-tier latency, the archive enumeration
threshold and the deduplication cooldown --- one at a time with everything else
at default.

\subsection{Reproducibility}

Every run is seeded and deterministic; a test asserts that two trials with the
same seed produce identical latencies and false-positive counts. Fourteen
simulated days take about two seconds of wall time. Every number in this paper is
generated from the experiment JSON by a script that emits the LaTeX macros and
table bodies the document includes, so a stale figure in the text is not possible.
